% 原书分章PDF第18—22页：第30.7节全部15段，所有引文与历史论证完整。
\section{消除、二元论、概率革命与统一化}

% S30.7.P01
因果性在人类行为和哲学中的显著作用，并未使这个概念免于在思想史进程中不断遭遇困境。十八世纪，大卫·休谟基于其概念经验主义，对因果性的（形而上学）地位提出了著名挑战。按照休谟的看法，因果性并不存在于实在之中；在我们的思想中，只有现象恒常而有规律的相继发生。由于自然习惯，它们在我们心中与物理必然性相联系，并被投射到实在之上，从而解释了我们对必然性的“直觉”。

% S30.7.P02
十九世纪，德国物理学家兼哲学家恩斯特·马赫，以及英国哲学家兼统计学家卡尔·皮尔逊，表现为著名的反因果论者。不过，最主要的批评者出现在二十世纪。正是前述伯特兰·罗素，于1913年在《亚里士多德学会会刊》（The Proceedings of the Aristotelian Society）发表题为《论原因的概念》（On the Notion of Cause）的札记，并说：“我相信，哲学家们所接受的因果律，是一个逝去年代的遗物；它像君主制一样继续存在，只因为人们错误地以为它没有害处。”这种对待因果性的态度并非独一无二。随着时间推移，许多概念与概念以及相关的哲学主题遭到否定，或至少被赋予一个特定的抽象或数学术语。这尤其适用于空间、时间、运动、意向性等概念，也适用于意义、心灵感、心灵、自由意志、意识与人格同一性等概念（Starmans，2011）。

% S30.7.P03
粗略地说，这是当代科学主义所致。一些科学家和自然主义哲学家认为，一切哲学问题最终都会被科学解开并解开。如果问题定义明确，它就会受到分析，并最终得到解决；如果定义不明确，它就会被认定为伪问题或毫无意义的问题，而被搁置一旁。我终将在时机成熟时消散（I will evaporate in the fullness of time）。哲学反思至多能形成某种前科学理论，或许具有一定解释力或实际用途，但最终会被真正的科学知识所取代。在这类哲学反思中发挥作用的概念与概念，通常也不得不消失。

% S30.7.P04
这种从科学及其语言中清除形而上学概念、常识概念、自然范畴和日常经验的倾向，在哲学家兼神经科学家保罗·丘奇兰（Paul Churchland）的观点中达到一个声名狼藉的高点。他要与一种有时被贬称为“民间心理学”的传统彻底决裂。人们试图用具有因果相关性的因素——如动机、意图、信念和义务——来理解、解释和预测自己与他人的行为。丘奇兰主张对这些命题态度采取激进的消除唯物主义，认为包括意识、人类心灵及其内部过程等概念的“民间心理学”完全不准确（Starmans，2011）。同样，他也抨击哲学家对语言及其被设定的、对于思维的关键意义的执着。按照丘奇兰的看法，神经科学的发展将消除这些“病症”；在他看来，它们与科学的相关程度，正如十八世纪燃素说与现代化学的相关程度，或中世纪巫术观念与当代心理学的相关程度一样。丘奇兰的极端看法绝非哲学家中的“共同意见”（communis opinion），但也并非一种新现象。它符合西方思想史上一段漫长而令人瞩目的发展：始于前苏格拉底哲学家，在今天的自然主义／物理主义认识论中达到顶峰。结果，当代世界观失去了许多直觉上的吸引力。一方面，是日常而熟悉的现象生活，产生经验（知觉、印象）、表征和观念，以及其中（被设定的）具体物质对象。另一方面，则是科学世界观及其抽象的、往往是数学的模型，它们表征一个被认为隐藏在这些经验背后、造成并解释这些经验的“真实”世界。

% S30.7.P05
这种二元论在德谟克利特（Demokritos）、留基伯（Leukippos）和伊壁鸠鲁的希腊原子论中已经广泛存在，并于十七世纪在笛卡尔和约翰·洛克的著作中以占主导地位的方式显现。他们区分了内在于物质的第一性质——不可分粒子的形状、大小、位置与运动——以及颜色、气味、声音等实际上只存在于人类心灵中的第二性质。两者之间的鸿沟，显然比以往任何时候都更大。我们所感知的这些现象，不仅不是科学理论的可靠基础，而且直觉概念与自然范畴，似乎也不符合科学语言及其术语体系所描述的“真实”世界的潜在机制、抽象原理和规律。

% S30.7.P06
事实上，因果性概念的整个演变，都必须置于这种张力之中来理解。首先，正是由于因果性在形而上学中的突出位置，它也成为消除企图的首要对象。当一种哲学立场遭到批评，或一种思辨理论被拆解时，与之相伴的因果性概念也立即受到判决。一个突出的例子是亚里士多德本人：他引入形式因与质料因，以支撑其连接形式与质料的质形论（hulemorphism）。此外，目的因对于他关于自然与宇宙中目的论／意向性的观点至关重要。

% S30.7.P07
随后，在十三世纪的经院哲学中，这个概念主要由阿奎那（Aquinos）重新解释和调整，以适合他把亚里士多德主义与基督教相结合的著名综合。但到了十七世纪，亚里士多德的物理学与形而上学遭到否定，受到培根、笛卡尔、霍布斯和惠更斯的批评。结果，只有动力因被认为可以接受。这一过程如此僵硬而仓促，以至于至今仍影响因果性问题。然而，培根也使因果性概念服从于其经验主义，假定存在某种可以直接从观察中推导出来的“致因形式”。笛卡尔则拥护一种严格的机械论哲学，其中，生产性的因果效应只能由粒子的碰撞、直接接触与瞬时效应来解释，不容许任何“超距作用”（actio in distans）；而后者在牛顿成功的反笛卡尔物理学中，当然具有关键地位。对于斯宾诺莎，因果性首先是一个逻辑概念。但培根、笛卡尔与斯宾诺莎的这些因果性观念，都不容许对实在作任何数学描述；在经典力学取得成功，以及拉普拉斯决定论随之兴起之后，它们便不再适用。

% S30.7.P08
休谟的怀疑论更多基于其概念经验主义，而非某种形而上学立场。在此之后，十八世纪主要是康德试图恢复这一概念，并把因果性视为一种先验综合知识和基本范畴，力求调和形而上学与自然科学。这构成观察与科学知识的一个条件。不过，随着二十世纪逻辑实证主义者的反康德哲学兴起，这种因果性观念也被认为存在问题；尤其因为量子力学取得了成功，其中不再把因果性视为康德意义上构成实在的基石。

% S30.7.P09
这个清单还可以很容易地继续扩展。历史有力地表明，因果性如何由于其形而上学根源，或与特定哲学立场的内在联系，而一次又一次地陷入困境。如前所述，这个概念不断被重新发明；但由于许多立场并未真正被证伪和取代，这些思想一直保留下来，演变为当今关于因果性的多元观点。因此，当代哲学家对因果性原则的关注，以及把握因果关系本质的努力，绝非仅仅出于历史兴趣。他们的思考与解释的某些方面，仍类似于现代精细的因果性进路，甚至类似于计算因果建模技术。

% S30.7.P10
几个例子是：反映休谟关于相继、邻近和恒常会合看法的规律性规则，反事实进路（也出自休谟），连续性与作用机制这一侧面（笛卡尔），以及干预解释（斯图亚特·密尔）。尤其应该肯定斯图亚特·密尔为弥合层次II与III之间的鸿沟而作出的努力。当十九世纪的哲学家开始思考实验研究时，他们需要一种新的、建设性的因果性概念，既能脱离旧形而上学，也能脱离拉普拉斯式决定论。在约翰·斯图亚特·密尔的著作中，可以看到把因果性锚定于更具实用性与实验性的语境中的尝试。他从确立因果关系或因果效应所必须认识、控制或操纵的全部情境，来理解因果性。他的《逻辑体系》（System of Logic，1843）包含著名的“密尔五法”，实际上是一本在“方法论手册”这一名称出现之前的方法论手册；作者试图弥合抽象认识论与实际科学实践之间的鸿沟，尤其是在因果性思考方面。亚里士多德及其《工具论》（Organon）、培根及其《新工具》（Novum Organum），还有阿尔诺与尼科尔及其《波尔—罗亚尔逻辑》（Logic of Port Royal）等杰出先驱，在这方面显然没有那么成功。不过，真正的进步还需要一场吸收概率论与统计学的概率革命。如今，几乎所有因果性进路都具有概率性质；没有概率革命，这个概念就绝不会在研究方法论中获得一席之地。

% S30.7.P11
十九世纪后期，对变异与变化的思考进一步加速，尤其得益于达尔文和华莱士的进化论。当然，当时这一理论根本没有进化的数学模型，也没有充分的遗传机制。正是弗朗西斯·高尔顿（Francis Galton），这位科学数学化的令人难以接受的倡导者（inadmissible advocate），认识到变异确实可以接受数学处理，无须带有异常、误差函数等贬义表现。高尔顿最先发展出“回归均值”（regression-to-the-mean）概念；他并不十分关心凯特勒（Quetelet）的平均值，而是关注偏离、特殊和个体，尽管他并不确切知道，长期的微小变化如何足以产生现象与进化的丰富多样性。

% S30.7.P12
另一方面，他的学生皮尔逊纠正了这一切：他不把变异定位于“误差”之中，而是定位于现象本身（由数据编码）之中，并区分不同的分布（类别）。他看到，许多现象并非正态，而是有偏的，可以用四个参数来分类：均值、标准差、偏度和曲率（curvature）。自然界的变异性表现为测量值组成的点云；皮尔逊寻求拟合得最好的模型，即最适合数据的函数。这与其说是在指示一个潜在的因果机制，不如说是对现象作一种简约描述。

% S30.7.P13
首先，他使概率在科学中获得完整地位，并在这样一个抽象层面看待世界：数据、数据变异、数据生成机制与分布参数，与其说是在表征或描绘一个被假定的物理世界，不如说是在编码和构建实在。概率分布及其相关参数，是科学的对象。世界观虽然失去了某些直接的可理解性，但矛盾的是，皮尔逊式的、作为统计分布的实在，却是“可观察的”、贴近数据的，尤其是可被注意到的。它描述实际数据，即所有可能测量值的集合中一个很大但有限的子集；所有可能的测量值，只有在一种“理想”情境下才能获得（Starmans，2011）。随着高尔顿和皮尔逊对变异与变化采取统计学进路，以及费希尔后来对进化论与孟德尔遗传学的综合，关于不确定性的思考发生了重大转变。在这里，我们只能说，费希尔的二元论立场有利于因果性思考，而皮尔逊的一元论立场则使因果性变得过时。迄今，因果性已经挺过了丘奇兰的消除论、冯·诺依曼遭到侵蚀的模型，以及 Anderson 对理论终结的思索；矛盾的是，这既不顾、又得益于它以某些至今仍流行且有必要的二元论形式出现。事实上，因果性的表现还要更好。作为一个对立面，甚至也许作为一剂解药，我们必须提及二十世纪末因果性思考中的计算转向。由此，科学家对这个从前棘手的概念有了更多把握。

% S30.7.P14
2012年，美国人工智能科学家 Judea Pearl 因其在概率网络与因果建模方面的工作，获得计算机科学最重要的科学奖项——图灵奖。2000年出版、2009年修订的《因果性：模型、推理与推断》（Causality: Models, Reasoning and Inference），可以视为现代因果性思考的一座里程碑（Pearl，2009a）。更进一步，Pearl 把因果性视为科学中一个具有统一作用的概念。事实上，他在导论中认为，“因果关系既是实在的基本构件，也是人类理解该实在的基本构件”。不可能把它还原为概率、统计学或算法。相反，Pearl 认为，“没有什么比把我们所有数学资源集中于概率与统计推断的流行做法，更能阻碍科学进步”；正是因为，按照他的看法，在这个传统中，因果概念及其术语被排斥在正式的科学语言之外。

% S30.7.P15
因此，他希望，在概率的二十世纪之后，二十一世纪将展现一个新的因果性时代。Pearl 追求统一化；在这些努力中，他把结构方程建模与概率网络相结合的框架发挥着关键作用。事实上，其方法的部分根源，可以在 Sewall Wright 的著作中找到。近一百年前，Wright 就已经建立了使用路径系数与线性方程组的方法，对不同内生变量与外生变量之间的关系、直接效应与间接效应等进行建模。Pearl 拒斥“对称的”数学语言，引入其著名的“do 算子”，以避免混杂及辛普森悖论等相关问题。毋庸置疑，这种颇为反统计学的立场，既丰富、又复杂化了人们对因果性的作用，以及它在统计学、数据科学与研究方法论中各种表现的看法。
